\section{引言}

张量链策略迭代（TTPI）利用低秩张量近似，使离散状态--动作空间中的混合控制问题变得可处理
\cite{shetty2024generalized,oseledets2011tensor}。其有利的复杂度依赖张量链表示中的秩，而这些秩并不是无序变量集合的内禀属性。它们由给定模态顺序诱导的一系列矩阵展开决定。因此，动作模态排序不是无关紧要的实现细节，而是 TTPI 中的一等算法变量。

最初的物理直觉很直接：通过动力学、奖励或价值备份强相互作用的动作变量，不应被分离到许多张量 cut 的两侧。然而，仅依靠物理信息构造排序目标并不等价于优化 TT 秩。加权线性排列刻画的是加权 cut profile 的总面积，而峰值显存和 TT-Cross 工作量往往由最坏 cut 决定。这一错配促使我们引入秩感知扩展：在保持混合执行器 block 连续的同时，优化 cut-level 瓶颈，而不仅仅优化物理邻近性。

因此，本文将方法定位为面向 TTPI 的秩感知、分块保留模态重排序。物理感知耦合作为先验使用；敏感度权重、奇异谱估计和实际 TT 秩诊断用于检验该先验是否与 TT 复杂度一致。当前仓库同步证据在 \path{notes/claims.md} 中以保守方式追踪：HardMove pilot 运行支持资源敏感性主张，但尚不支持排序能提升最终控制性能的主张。
